\section{Instantiation}\label{sec:agg-creds-instantiation}

\paragraph{Notation.} We use the bilinear groups and conventions of the Preliminaries. In this chapter, elements of $\group{1}$ are
written using lower case letters, elements of $\group{2}$ carry a tilde, and
elements of $\Zp$ are written using lower case Greek letters or otherwise
specified explicitly. The exceptions to this are $c$, which is a challenge in
$\Zp$, and $\signature$ and $\SPSsig$, which are signatures consisting of
several group elements. We still use a tilde on the elements of the signatures
to signal which groups they belong to.

We use hash functions $\hashone \colon \{0,1\}^* \rightarrow \group{1}$ and
$\hashtwo \colon \{0,1\}^* \rightarrow \group{1}$, and hash functions $\hash
\colon \{0,1\}^* \rightarrow \Zp$ for the challenges of the proofs below.
Each hash function prefixes its input with its own fixed string, its domain
separation tag, so that they behave as independent random oracles.

\subsection{Instantiating $\ATS$}\label{sec:agg-creds-tag-scheme}

Our tag-based aggregate signature scheme borrows its structure from the BLS
signature scheme of Boneh, Lynn and Shacham and its aggregate form, in which a
signature is the hash of the message raised to the secret key and signatures on
different messages aggregate by multiplication~\cite{JC:BonLynSha04,EC:BGLS03}.
Each of our signatures keeps this term and adds the hash of a tag raised to a
fresh exponent, so that only signatures under the same tag aggregate and
signatures can be rerandomised.

\begin{outline}
  \item why this structure: signatures aggregate by multiplication and verify with one pairing equation, so aggregating costs one group multiplication per signature and the aggregate is as short as one signature
\end{outline}
Both hash functions are modelled as random oracles, and the security proof
embeds its challenge in the hashes of the attacked message and tag. Our scheme is comprised of the following algorithms:

\begin{protocolbox}{The tag-based aggregate signature scheme}{ats}

  \begin{description}

    \item[$(\sk, \pk) \sample \ATSgen(\secparam)$ :] On input the security
parameter, do the following:

      \begin{enumerate}

	\item Sample $\sk \sample \Zp$.

	\item Return $(\sk, \pk = \ggen^\sk)$.

      \end{enumerate}

    \item[$\signature \sample \ATSsign(\sk, \Tag, \msg)$ :] On input secret key
$\sk$ and pair of tag and message $(\Tag, \msg)\in \{0,1\}^*\times\{0,1\}^*$,
do:

      \begin{enumerate}

	\item Sample $\upsilon \sample \Zp$.

	\item Return $\signature = (\ggen^\upsilon, \hashone(\msg)^\sk \cdot
\hashtwo(\Tag)^\upsilon)$.

      \end{enumerate}

    \item[$\signature' \sample \ATSrandomise(\Tag, \signature) $ :] On input tag
$\Tag$, signature $\signature$, do:

      \begin{enumerate}

	\item Parse $\signature = (\signatureone, \signaturetwo)$.

	\item Sample $\upsilon' \sample \Zp$.

	\item Return $\signature' = (\signatureone \cdot  \ggen^{\upsilon'},
\signaturetwo \cdot \hashtwo(\Tag)^{\upsilon'})$.

      \end{enumerate}

    \item[$ \{0,1\} \leftarrow \ATSverify(\pk, \Tag, \msg, \signature) $ :] On
input public key $\pk$, tag $\Tag \in \{0,1\}^*$, message $\msg \in \{0,1\}^*$,
and signature $\signature$, do:

      \begin{enumerate}

	\item Parse $\signature = (\signatureone, \signaturetwo)$.

	\item Return $\pairing(\signaturetwo,\ggen) \checkcheck
\pairing(\hashone(\msg), \pk) \cdot \pairing(\hashtwo(\Tag), \signatureone)$.

      \end{enumerate}

    \item[$\signature \leftarrow \ATSaggregate((\signature_i)_{i \in
\ISSUERS})$:] On input signatures $(\signature_i)_{i \in \ISSUERS}$, for set
$\ISSUERS \subseteq [1, n]$, do:

      \begin{enumerate}

	\item Parse $\signature_i = (\signatureone_i, \signaturetwo_i)$ for each
$i \in \ISSUERS$.

	\item Return $\signature = (\prod_{i\in\ISSUERS} \signatureone_i, \
\prod_{i\in\ISSUERS} \signaturetwo_i)$.

      \end{enumerate}

    \item[$\{0,1\} \leftarrow \ATSverifyagg((\pk_i)_{i\in \ISSUERS},
\Tag,(\msg_i)_{i\in \ISSUERS},\signature) $ :] On input public keys
$(\pk_i)_{i\in \ISSUERS}$, tag $\Tag$, messages $(\msg_i)_{i\in \ISSUERS}$, and
aggregate signature $\signature$, the aggregate verification algorithm outputs
the result of the following check:

      \begin{enumerate}

	\item Parse $\signature = (\signatureone, \signaturetwo)$.

	\item Return $\pairing(\signaturetwo, \ggen) \checkcheck \pairing
(\hashtwo(\Tag), \signatureone) \cdot \prod_{i\in\ISSUERS}
\pairing(\hashone(\msg_i), \pk_i)$.

      \end{enumerate}

  \end{description} \end{protocolbox}

This signature is existentially unforgeable under tag-based aggregation in the
random oracle model under the Bilinear Computational Diffie-Hellman assumption
and a modified variant of the same assumption.

  \begin{assumptionbox}{Bilinear Computational Diffie-Hellman (BCDH)}{bcdh}

    The BCDH assumption holds if, given $(\gen, \gen^\gamma, \gen^\varepsilon,
\gen^\chi, \ggen, \ggen^\gamma, \ggen^\varepsilon, \ggen^\chi)$ for uniformly
random $\gamma, \varepsilon, \chi \sample \Zp$, no PPT adversary can output
$\pairing(\gen, \ggen)^{\gamma\varepsilon\chi}$ with non-negligible probability.

  \end{assumptionbox}

  \begin{assumptionbox}{Modified Bilinear Computational Diffie-Hellman
(m-BCDH)}{mbcdh}

    The m-BCDH assumption holds if, given $(\gen, \gen^{1/\varepsilon}, \ggen,
\ggen^\gamma, \ggen^\varepsilon)$ for uniformly random $\gamma, \varepsilon
\sample \Zp$, no PPT adversary can output $\pairing(\gen,
\ggen)^{\gamma\varepsilon}$ with non-negligible probability.

  \end{assumptionbox}

  \begin{theorem}\label{thm:agg-creds-thm-1} If the BCDH and m-BCDH assumptions
hold, then our tag-based aggregate signature scheme is existentially unforgeable
under tag-based aggregation in the random oracle model.  \end{theorem}

The proof separates two kinds of forgery. BCDH covers a forgery on a message
never signed under the honest key, and m-BCDH covers a signed message reused
under a new tag.
\begin{outline}
  \item neither assumption is known to imply the other, so both are needed
\end{outline}
Both assumptions hold in the GGM, the model in which the full credential scheme
is proven secure.

\subsection{Instantiating $\AAC$ with $\AGHO$ and $\ATS$}\label{sec:agg-creds-aac-inst}

\paragraph{Overview.} In this instantiation, each user has a pair of secret and public keys $(\sk, \pk
= \gen^\sk)$. The digital signature $\DS$ uses the AGHO
structure-preserving signature scheme~\cite{C:AGHO11}. Every user is assigned a
unique tag $\Tag \in \Zp$. Since structure-preserving signatures sign vectors
of group elements rather than scalars, we take $\tid = \hashtwo(\Tag)$ and it is
the pair $(\pk, \tid)$ that the registration authority signs, with $\pk$ the
user's public key.  Every other computation in the protocol uses $\Tag$
directly. The aggregate tag-based signature $\ATS$ is instantiated with our
construction. The protocol instantiation is described in full in
Protocol~\ref{prot:paac}.

\paragraph{The AGHO structure-preserving signature scheme.} The registration authority's key pair is for the structure-preserving signature
scheme, since messages and signatures consist of elements of the same group as
everything else in the protocol. This lets users prove knowledge of a valid
registration signature with an efficient sigma protocol over discrete
logarithms, rather than requiring a general purpose zero knowledge proof.

The credential scheme is then only proven secure in the GGM, since the $\AGHO$
SPS scheme is only proven secure there.

During registration, the registration authority signs the user's public key
together with the group element $\tid = \hashtwo(\Tag)$, rather than $\Tag$
itself. This serves the same purpose as above, and lets us extract $\Tag$ from
$\tid$ by programming the random oracle. All other computations in the protocol,
including the tag-based aggregate signature scheme, use the tag directly.

\begin{outline}
  \item why hash rather than $\gen^\Tag$: $\tid$ then has no algebraic relation to the other exponents in the protocol
\end{outline}
The aggregate signature uses the same element $\hashtwo(\Tag)$, so the
registration signature and the aggregate signature share one group element, and
the presentation only shows that they share it.
Since the registration authority only issues signatures on well-formed
messages, pairs $(\pk, \hashtwo(\Tag))$ for $\Tag \in \Zp$, the user does not
need to prove knowledge of the preimage such that $\tid = \hashtwo(\Tag)$,
which already saves a proof of a hash preimage.

At presentation, users raise every element that would identify the
rerandomised signatures to a fresh $\delta$, including $\gen$ and the attribute
hashes. The Sigma protocol proves $\rel_3$: that one $\delta$ is used
throughout, and knowledge of $\sk$. The challenge hashes every group element of
the presentation and the verifier's message, so nothing can be changed once the
challenge is fixed.

In the concrete instantiation, the user randomises both signatures, to
$(\signatureone', \signaturetwo')$ and $(\widetilde{r}', \widetilde{s}', z')$. This means that these values can be sent in the clear within
$\algoproof$, as they reveal nothing about which original signature they came
from. This means that the signature verification equations can be used
directly. $\algoproof$ then captures the full presentation relation of the
construction, with signature verification performed directly instead of proven,
and the discrete-log equalities enforcing consistency through $\eta$. This
allows us to build a concretely efficient proof rather than requiring
zero-knowledge proofs over pairing equations and elements of $\group{T}$.

If a user is able to produce a valid $\eta$, we can extract a witness for the
presentation relation of the construction. Under the DDH assumption in
$\group{1}$, the resulting presentation is simulatable: it leaks no information
about the witness beyond the disclosed attribute values.

\begin{lemma}\label{lem:agg-creds-sigma}

  A witness for the presentation relation of the construction can be extracted
from any presentation that verifies. If the DDH assumption holds in
$\group{1}$, presentations can be simulated without a witness.

\end{lemma}

\Cref{lem:agg-creds-sigma} gives extractability and simulatability.
\Cref{thm:agg-creds-thm-2} needs simulation extractability, which additionally
allows extraction after the adversary has observed simulated proofs of its own
choosing. Our simulator does both within the same experiment, so this is the
property we rely on.  The proof in Protocol~\ref{prot:paac} is a sigma protocol with the
Fiat-Shamir heuristic applied~\cite{C:FiaSha86}. This combination satisfies
simulation extractability in the random oracle model~\cite{INDOCRYPT:FKMV12}.

At registration and issuance, users prove knowledge of $\sk$ with a Sigma
protocol on a nonce chosen by the registration authority or the issuer, so proofs
cannot be replayed. Issuers check the registration signature before signing, so
they keep no record of whom they have issued to.

%%%% anon creds protocol

\begin{protocolbox}{The aggregate anonymous credentials protocol}{paac}

  \begin{description}

    \item[$\crs \sample \AACgen(\secparam, n)$ :] On input security parameter
$\secparam$ and number of issuers $n$, do the following:

      \begin{enumerate}

	\item Sample $(\rask,\rapk) \sample \SPSgen(\secparam,2)$ for the
$\AGHO$ SPS scheme, with $\rapk = (\widetilde{u}_1, \widetilde{u}_2, v)$.

	\item For each issuer $i \in [1,n]$, sample $(\isk{i},\ipk{i}) \sample
\ATSgen(\secparam)$, so $\ipk{i} = \ggen^{\isk{i}}$.

	\item Return $\crs \leftarrow
(\gen,\ggen,\hashone,\hashtwo,\hash,\rapk,\{\ipk{i}\}_i)$.

      \end{enumerate}

    \item[$(\sk,\regcred) \sample \AACregister(\crs)$ :] Between the user and
the RA, do:

      \begin{enumerate}

	\item The user samples $\sk \sample \Zp$ and sets $\pk \leftarrow
\gen^\sk$.

	\item The RA samples a fresh nonce $\rand \sample \Zp$ and sends it to
the user.

	\item The user samples $\varrho \sample \Zp$, computes $a = \gen^\varrho$,
$c = \hash(\gen, \pk, a, \rand)$ and $\zeta = \varrho + c \cdot \sk$, and sends
$(\pk, a, \zeta)$.

	\item The RA recomputes $c$ and exits unless $\pk \neq 1$ and $\gen^\zeta
= a \cdot \pk^c$.

	\item Otherwise, the RA samples a fresh $\Tag \sample \Zp$ it has not
assigned before, sets $\tid \leftarrow \hashtwo(\Tag)$, and computes $\SPSsig
\sample \SPSsign(\rask,(\pk,\tid))$.

	\item Return $\sk$ and $\regcred \leftarrow (\Tag,\pk,\SPSsig)$ to the
user.

      \end{enumerate}

    \item[$\cred \sample \AACissue(\crs,\isk{},\regcred,\att{})$ :] Between
the user, on input $(\sk, \regcred)$, and issuer $i$, on input $\isk{i}$, for
the user's value $\atv{i}$ of attribute $\att{i}$, do:

      \begin{enumerate}

	\item The user sends $\regcred = (\Tag,\pk,\SPSsig)$. The issuer
computes $\tid \leftarrow \hashtwo(\Tag)$ and exits if
$\SPSverify(\rapk,(\pk,\tid),\SPSsig) \neq 1$.

	\item The issuer sends a fresh nonce $\rand$, and the user proves
knowledge of $\sk$ as in $\AACregister$. If the proof fails, the issuer exits.

	\item Otherwise, the issuer computes $\signature_i \sample
\ATSsign(\isk{i},\Tag,\atv{i})$, and the user stores $\cred_i \leftarrow
(\atv{i},\signature_i)$.

      \end{enumerate}

    \item[$\aggcred \leftarrow \AACaggregate(\crs,(\cred_i)_{i \in \ISSUERS})$ :]
On input the credentials $(\cred_i)_{i \in \ISSUERS}$ to be disclosed, for
$\ISSUERS \subseteq [1,n]$, all issued under the same tag $\Tag$, do:

      \begin{enumerate}

	\item Parse each $\cred_i = (\atv{i},\signature_i)$ for $i \in \ISSUERS$.

	\item Compute $\signature \leftarrow
\ATSaggregate((\signature_i)_{i\in\ISSUERS})$.

	\item Return $\aggcred \leftarrow
(\Tag,\{(i, \atv{i})\}_{i\in\ISSUERS},\signature)$.

      \end{enumerate}

    \item[$\algoproof \sample \AACpresent(\crs,\sk,\regcred,\aggcred,\msg)$ :]
On input the user's $\sk$ and $\regcred$, $\aggcred$, and message $\msg$
from the verifier, do:

      \begin{enumerate}

	\item Parse $\regcred = (\Tag,\pk,\SPSsig)$ and $\aggcred =
(\Tag,\{(i, \atv{i})\}_{i \in \ISSUERS},\signature)$, and set $\tid =
\hashtwo(\Tag)$.

	\item Randomise the signatures: $(\signatureone', \signaturetwo')
\sample \ATSrandomise(\Tag,\signature)$ and $(\widetilde{r}', \widetilde{s}',
z') \sample \SPSrandomise(\SPSsig)$.

	\item Sample $\delta \sample \Zp^*$ and compute
          %
	  \begin{gather*}
          %
	    \bar{\gen} = \gen^\delta, \quad \bar{v} = v^\delta, \quad \bar{z} =
{z'}^\delta, \quad \bar{\signaturetwo} = {\signaturetwo'}^\delta, \\
\bar{\tid} = \tid^\delta, \quad \pk' = \pk^\delta, \quad H_i =
\hashone(\atv{i})^\delta \ \ \forall i \in \ISSUERS.
          %
	  \end{gather*}

	\item Set the statement and witness
          %
	  \begin{equation*}\begin{gathered}
          %
	    \inst_3 = (\bar{\gen}, \bar{v}, \pk', \{i, \atv{i}, H_i\}_{i\in\ISSUERS},
\widetilde{r}', \widetilde{s}', \bar{z}, \signatureone', \bar{\signaturetwo},
\bar{\tid}), \quad \witness_3 = (\delta, \sk), \\
          %
	    \rel_3 = \{\ \bar{\gen} = \gen^\delta \ \wedge\ \bar{v} = v^\delta \
\wedge\ H_i = \hashone(\atv{i})^\delta\ \ \forall i \in \ISSUERS \ \wedge\ \pk'
= \bar{\gen}^\sk\ \}.
          %
	\end{gathered}\end{equation*}

      \item Prove $\rel_3$: sample $\alpha, \beta \sample \Zp$, compute
          %
	  \begin{equation*}
          %
	    (f, h, \{k_i\}_{i \in \ISSUERS}, w) = (\gen^\alpha, v^\alpha,
\{\hashone(\atv{i})^\alpha\}_{i \in \ISSUERS}, \bar{\gen}^\beta),
          %
	  \end{equation*}
          %
	the challenge $c \leftarrow \hash(\gen, v, \inst_3, f, h, \{k_i\}_{i \in
\ISSUERS}, w, \msg)$, and the responses $\pi_1 = \alpha + c\delta$ and $\pi_2 =
\beta + c\sk$, and set $\eta = (f, h, \{k_i\}_{i \in \ISSUERS}, w, \pi_1,
\pi_2)$.

      \item Return $\algoproof = (\eta, \inst_3)$.

    \end{enumerate}

  \item[$\{0,1\} \leftarrow \AACverify(\crs,\msg,\algoproof)$ :] On input
$(\msg,\algoproof)$ from a verifier, parse $\algoproof = (\eta, \inst_3)$ with
$\inst_3 = (\bar{\gen}, \bar{v}, \pk', \{i, \atv{i}, H_i\}_{i\in\ISSUERS},
\widetilde{r}', \widetilde{s}', \bar{z}, \signatureone', \bar{\signaturetwo},
\bar{\tid})$ and $\eta = (f, h, \{k_i\}_{i \in \ISSUERS}, w, \pi_1, \pi_2)$, and do:

    \begin{enumerate}

      \item If $\bar{\gen} = 1$, return $0$.

      \item Compute $c \leftarrow \hash(\gen, v, \inst_3, f, h, \{k_i\}_{i \in
\ISSUERS}, w, \msg)$.

      \item Return the output of the following checks:
       %
	\begin{gather*}
		%
	  \gen^{\pi_1} \checkcheck f\cdot\bar{\gen}^c, \quad v^{\pi_1}
\checkcheck h\cdot\bar{v}^c, \quad \hashone(\atv{i})^{\pi_1} \checkcheck
k_i\cdot H_i^c \ \ \forall i \in \ISSUERS, \quad \bar{\gen}^{\pi_2}
\checkcheck w\cdot\pk'^c, \\
		%
	  \pairing(\bar{v}, \widetilde{r}')\checkcheck \pairing(\bar{\gen},
\widetilde{s}'), \\
		%
	  \pairing(\bar{\signaturetwo}, \ggen) \checkcheck \prod_{i\in
\ISSUERS}\pairing(H_i, \ipk{i})\cdot \pairing(\bar{\tid}, \signatureone'), \\
          	%
	  \pairing(\bar{\gen}, \ggen) \checkcheck \pairing(\bar{z},
\widetilde{r}')\cdot \pairing(\pk', \widetilde{u}_1)\cdot \pairing(\bar{\tid},
\widetilde{u}_2).
	%
	\end{gather*}
        %
    \end{enumerate}

\end{description}

\end{protocolbox}

\begin{theorem}\label{thm:agg-creds-thm-2}

    In the GGM and the random oracle model, the aggregate anonymous credentials
protocol $\AAC$ of Protocol~\ref{prot:paac} securely realises ideal functionality
$\idfu{\AAC}$.

\end{theorem}

